\section{Stabilization inside the exponent box}
\label{sec:stabilization}

We now start with a monomial ideal containing the prescribed powers.
The aim of this section is to impose stability inside the finite exponent
box while preserving the Hilbert function and weakly increasing every
graded Betti number of the quotient. This is the pure-power
deformation strategy of Mermin and Murai \cite[Section~3]{MM}.

Let
\[
 S=\C[x_1,\ldots,x_n],\qquad
 2\le a_1\le\cdots\le a_n,\qquad
 P=(x_1^{a_1},\ldots,x_n^{a_n}),\qquad b_i=a_i-1.
\]
For $1\le r\le n$ and $d\in\Z$, put
\[
 \mathcal B_r(d)
 =\{u\in\Z_{\ge0}^r: |u|=d,\ 0\le u_i\le b_i
                          \text{ for }1\le i\le r\}.
\]
In particular, this set is empty when $d<0$.

\begin{definition}
\label{stab:definition}
A subset $X\subseteq\mathcal B_r(d)$ is \emph{stable} if
\[
 u\in X,\quad h<\ell,\quad
 u+e_h-e_\ell\in\mathcal B_r(d)
 \quad\Longrightarrow\quad u+e_h-e_\ell\in X.
\]
A monomial ideal containing $(x_1^{a_1},\ldots,x_r^{a_r})$ is
\emph{stable in the box} if its included exponent vectors in
$\mathcal B_r(d)$ form a stable set for every $d$.
This definition imposes no transfer condition on monomials outside the box.
\end{definition}

The modification used below acts on one pair of variables.  Fix
$h<\ell$, write $\xi=x_h$ and $q=x_\ell$, and choose a primitive
$a_\ell$-th root of unity $\zeta\in\C$.  Define a graded
$\C$-linear map $D:S\to S$ on monomials by
\begin{equation}
\label{stab:distraction-formula}
 D(x^u)=\left(\prod_{i\ne\ell}x_i^{u_i}\right)
             \prod_{s=0}^{u_\ell-1}(q-\zeta^s\xi),
 \qquad u\in\Z_{\ge0}^n,
\end{equation}
with the empty product equal to $1$.
The resolution argument below is the nested-factor mechanism used for
distractions in~\cite[Theorem~2.19 and Corollary~2.20]{BCR}.
We give the argument directly because the
root-of-unity factors here are periodic, rather than eventually constant.

\begin{lemma}
\label{stab:distraction}
For every monomial ideal $K\subseteq S$, its vector-space image $D(K)$
is a homogeneous ideal.  Moreover,
\begin{gather*}
 \dim_\C(S/D(K))_d=\dim_\C(S/K)_d
 \quad(d\in\Z),\\
 \beta^S_{p,j}(S/D(K))=\beta^S_{p,j}(S/K)
 \quad(p\ge0,\ j\in\Z).
\end{gather*}
The map also satisfies $D(P)=P$.
\end{lemma}

\begin{proof}
The expansion of $D(x^u)$ has coefficient $1$ at $x^u$, and every
other monomial is lexicographically larger than $x^u$.  Thus $D$ is
unitriangular on each finite-dimensional graded piece and is invertible.
Multiplication by a variable other than $q$ commutes with $D$, while
\begin{equation}
\label{stab:multiplication}
 qD(x^u)=D(qx^u)+\zeta^{u_\ell}\xi D(x^u).
\end{equation}
Since $K$ is spanned by monomials and is an ideal, these identities show
that $D(K)$ is an ideal.  The invertibility and degree preservation of
$D$ give the Hilbert-function equality.

To check preservation of $P$, first observe that
\[
 \prod_{s=0}^{a_\ell-1}(q-\zeta^s\xi)
       =q^{a_\ell}-\xi^{a_\ell}\in P.
\]
Here the second term belongs to $P$ because $a_h\le a_\ell$.
If a monomial is divisible by $q^{a_\ell}$, its image under $D$ is
divisible by this displayed polynomial.  If it is divisible by another
generator of $P$, that generator remains a factor of its image.
Consequently $D(P)\subseteq P$.  In each degree this inclusion is an
inclusion between vector spaces of the same finite dimension, so it is
an equality.

It remains to compare Betti numbers.  If $K=S$, both quotients are zero.
Otherwise, choose a minimal exponent-multigraded free resolution
$F_\bullet\to S/K$ with $F_0=S$, as in
Lemma~\ref{filt:minimal}.  A nonzero differential entry from a generator
of multidegree $u$ to one of multidegree $v$ has the form
$c x^{u-v}$, where $u\ge v$ coordinatewise.  Keeping the ordinary
graded free modules, replace this entry by
\begin{equation}
\label{stab:ratio-entry}
 c\frac{D(x^u)}{D(x^v)}.
\end{equation}
The ratio is a polynomial: all unchanged monomial factors are nested,
as are the products in~\eqref{stab:distraction-formula}.

For each free summand whose generator has multidegree $v$, define a
linear change on its coefficient polynomials by
\begin{equation}
\label{stab:coefficient-change}
 C_v(x^{w-v})=\frac{D(x^w)}{D(x^v)},\qquad w\ge v.
\end{equation}
More explicitly, the right-hand side is
\[
 \left(\prod_{i\ne\ell}x_i^{w_i-v_i}\right)
       \prod_{s=v_\ell}^{w_\ell-1}(q-\zeta^s\xi).
\]
It has the original monomial $x^{w-v}$ with coefficient $1$ and only
lexicographically larger additional terms.  Hence $C_v$ is an
invertible graded vector-space map.  These maps need not be
$S$-linear; the new differential defined by~\eqref{stab:ratio-entry}
is $S$-linear.

The changes on the free summands intertwine the old and new
differentials.  Indeed, for a source coefficient $x^{w-u}$ and a
nonzero entry $c x^{u-v}$, the required equality is precisely
\[
 \frac{D(x^w)}{D(x^u)}
 \left(c\frac{D(x^u)}{D(x^v)}\right)
      =c\frac{D(x^w)}{D(x^v)}.
\]
Thus the new maps form a complex that is exact in every positive
homological degree.  At $F_0=S$ the change is $C_0=D$, so its image
from $F_1$ is $D(K)$.  We have obtained a free resolution of $S/D(K)$
with the same ordinary graded free modules as $F_\bullet$.
Every replaced nonzero entry still has its original positive degree,
so this resolution is minimal.  Its free ranks in each degree give
the asserted equality of all graded Betti numbers.
\end{proof}

\begin{proposition}
\label{stab:main}
Let $S=\C[x_1,\ldots,x_n]$, let $2\le a_1\le\cdots\le a_n$,
and put $P=(x_1^{a_1},\ldots,x_n^{a_n})$.
For every monomial ideal $K\supseteq P$, there is a monomial ideal
$H\supseteq P$ stable in the box such that
\[
 \dim_\C(S/H)_d=\dim_\C(S/K)_d\quad(d\in\Z)
\]
and
\[
 \beta^S_{p,j}(S/K)\le\beta^S_{p,j}(S/H)
 \qquad(p\ge0,\ j\in\Z).
\]
\end{proposition}

\begin{proof}
If $K$ is stable in the box, take $H=K$.  Otherwise choose a failed
transfer from $q=x_\ell$ to $\xi=x_h$, with $h<\ell$, and form
the map $D$ above.  Replace $K$ by
\[
 K'=\operatorname{in}_{\mathrm{lex}}D(K),
\]
where the leading monomial is the lexicographically largest one.
By Lemma~\ref{stab:distraction}, $P\subseteq D(K)$, and hence
$P\subseteq K'$.  That lemma and Proposition~\ref{filt:comparison}
show that the Hilbert function is preserved and
\[
 \beta^S_{p,j}(S/K)\le\beta^S_{p,j}(S/K')
 \qquad(p\ge0,\ j\in\Z).
\]
We prove that repeated replacements terminate by measuring progress
in every lexicographic prefix.

Fix a degree $d$ and a prefix $E$ of the degree-$d$ monomials, listed
from greatest to least.  Let $\pi_E$ be the projection onto their span.
Gaussian elimination with this ordered list of columns gives
\[
 \#\{m\in E:m\in K'\}
       =\dim_\C\pi_E(D(K)_d).
\]
For the monomials $m\in K\cap E$, the vectors $\pi_E(D(m))$
are linearly independent: their coordinates in their original
monomial positions form a triangular matrix with diagonal entries $1$.
Consequently
\begin{equation}
\label{stab:prefix-growth}
 \#(K'\cap E)\ge\#(K\cap E)
\end{equation}
for every such prefix in every degree.

There is a strict inequality for at least one prefix.  Choose a
monomial $x^u\in K$ inside the box witnessing the failed transfer, so
\[
 m'=\frac{\xi}{q}x^u
 \quad\text{lies inside the box but does not belong to }K.
\]
Its existence implies $1\le u_\ell<a_\ell$.  Let $E$ be the
degree-$|u|$ prefix ending at $m'$.  Every monomial in $K\cap E$
is strictly larger than $m'$, so its $D$-image has zero coefficient
at $m'$.  On the other hand, the coefficient at $m'$ of $D(x^u)$ is
\[
 -\sum_{s=0}^{u_\ell-1}\zeta^s
       =-\frac{1-\zeta^{u_\ell}}{1-\zeta}\ne0.
\]
Thus $\pi_E(D(x^u))$ is independent of the projected vectors already
used for $K\cap E$, proving strictness in~\eqref{stab:prefix-growth}.

Every monomial ideal containing $P$ is determined by the monomials
it includes in the finite box $0\le u_i<a_i$.  There are therefore
only finitely many such ideals.  The nondecreasing prefix counts,
with at least one strict increase at every replacement, prevent a
repetition.  Continuing whenever stability fails must consequently
stop at an ideal $H$ stable in the box.  The Hilbert-function equalities
and all Betti inequalities persist along this finite sequence.
\end{proof}
